\documentclass[11pt]{article}
\usepackage{amsmath,amssymb,amsthm,geometry}\geometry{margin=1in}
\newtheorem{proposition}{Proposition}
\title{Marked pure boundary protocols and response-domain reindexing}
\date{17 November 2026}\begin{document}\maketitle
\paragraph{Scope.}
This checks the common-base obligation of \texttt{Schema\_Relative\_Crossing\_Payload.tex}, using the independently established reserve-saturated atomic polynomial construction. It is a geometry and response-domain check, not a replacement proof of the paid conditional-family or balanced-crossing equations. The distinction between structural invocation slots and actual primitive observations is retained throughout.

\section{Actual marked base}
Take the realized reserve-saturated intrinsic protocol category: vertices, entry symbols with selected-entry incidence, and atomic edges over one common ticket; uniform maps on these sorts commute with incidence and preserve the declared actual operational observations, primitive supports, sponsor/predecessor fields and atomic sorts. Footprint eligibility is a positive requirement on OLD registry entries; new support has the separate current-event stamp. Maps satisfy source-eligibility implies target-eligibility. They need not reflect eligibility. All this is law-free. Pending entry symbols are not receiver initializer witnesses.

Add a weakening-invariant vertex predicate $c_G$. Restrict maps by
\[
                   c_H(fv)=c_G(v).
\]
Define the predicate of a polynomial vertex by
\[
 c_{P_CG}(u(v_1,\ldots,v_n))=0,\qquad
 c_{P_CG}(d(a;v_1,\ldots,v_n))=\bigwedge_i c_G(v_i).
\]
The empty conjunction is true. The predicate is not an extra strict requirement to preserve every latent parent phase separately.

\begin{proposition}
The reserve-saturated polynomial, its unit and multiplication restrict to this marked category. Their graph pullback and cartesian naturality statements remain valid. The pure discrete branch contains every Set function and every original arity-varying content body.
\end{proposition}
\begin{proof}
Polynomial maps preserve and reflect the displayed conjunction because their leaf maps do. The unit inserts a completed unary parent. The nested predicate for multiplication is the conjunction of all parent-completion flags and all final leaf predicates. The flattened predicate has exactly that value, including empty groups: an unfinished empty group still prevents parent completion. Hence multiplication is an exact-mark map. Its existing incidence, entry and footprint calculations are unchanged.

A pullback pair has equal marks, since both maps are exact-mark. Give it that mark. The previous unique tuple/edge reconstructions respect these marks; for a multiplication lift the supplied nested item fixes its parent flags and the reconstructed leaves have the matching marks. Thus the same pullback proofs restrict.

In $jS$, every value vertex is closed and strict operational output and primitive support are empty. Entry symbols represent the same pure values. A function of sets preserves these fields without being injective. A content body $a\mapsto d(c;b_1,\ldots,b_n)$ has closed target for every $n$, including zero; the tuple coordinates are structural, not a claimed bijection from one actual actor to $n$ actors. Its entry map is the corresponding generated tuple of formal entry symbols. Full collection shapes, including unused shapes and their coefficient arrows, remain in the content grade, not in a receiver law.\end{proof}

\section{Entry and passive response domains are pulled back, not guessed}
An entry domain is a specified collection of typed entry inputs (including their entry incidence), and a passive-response domain at state $x$ is a specified collection of boundary paths with actual responses. It is NOT the set of all paths with no visible output growth. For a pure input map $f:A'\to A$, use its entry-sort map and free-path map, together with its complete-signature map $f_\#$.

At pulled-back state $(a',x)$, define each accepted entry or response domain to be the inverse image of the corresponding domain at $(f(a'),x)$. The response is the original response, with the source postprofile $a''$ retained in the pulled-back poststate $(a'',x'')$. This is a definition on every raw atom, not a restriction to positive support. In particular an undefined original response remains undefined at all its preimages. No path lift, selected law or normalization is needed.

\begin{proposition}
These inverse-image domains preserve partial response composition and all-raw response certificates, and compose under pure input maps. They preserve reserve locality provided the pure maps are uniform intrinsic maps of the stated saturated interface. Output postcomposition leaves the response domains unchanged.
\end{proposition}
\begin{proof}
A composable admitted source path pair maps to the corresponding composable target pair; the original response equation supplies the result and the retained source endpoints type its pullback. Inverse images commute with composition of maps. Entry selected-state incidence is the same argument on the entry sort. Every complete source row is forwarded as a whole signature by $f_\#$; none of its zero atoms is removed. Exact closed marks ensure that completion of the original state at $f(a')$ implies closed input $a'$. Output postcomposition preserves closed output and its incidence equations.

For reserve locality, weakening changes only the common ticket tag of an intrinsic item. Uniformity means $f(\operatorname{rep}(t,x))=\operatorname{rep}(t,fx)$, so the same intrinsic local consumed tests and their boundary-wire pullbacks are used at all enlargements. Whole raw/law/factor reserve is still forwarded. Eligibility implication guarantees that every source representative has a target representative; target extra representatives need not be reflected. This argument would fail for a map selecting different intrinsic profiles according to unrelated reserve contents. Such maps are excluded, not declared harmless by postcomposition notation.\end{proof}

This does not assert that a partial receiver is receptive to every pure boundary. Its pulled-back domain may be empty. It also does not convert a virtual birth into a primitive passive response: receiver entry schemas realize virtual invocations later, with their own primitive variables, selected states and material factors. An empty schema realizes no primitive request, while a genuine unit-cost primitive remains a real occurrence.

\section{Two exact limits on stronger interpretations}
\paragraph{Footprint reflection fails for multiplication.}
Let an outer atomic edge add a new EMPTY inner group with sponsor $s$. Its nested eligibility includes the fact that $s$ is an active sponsor. Flattening produces no new leaf request and can drop that empty-wrapper requirement. A flat representative can therefore exist at a ticket lacking $s$, although no nested representative does. Multiplication still defines a map at every ticket admitting its source. Consequently exact equality of eligible-ticket domains, unlike the required implication, is not closed under multiplication. The cartesian lifting theorem recovers the nested sponsor from a SUPPLIED nested target and does not contradict this example.

\paragraph{Virtual stationarity is not primitive passivity.}
In the same example the outer virtual injection grows from zero to one group, whereas the flattened leaf injection stays empty. If ``passive'' were an exact reflected predicate meaning ``no new virtual slots'', multiplication would fail to reflect it. The primitive passive-response sort and its independently supplied receiver domain must instead be retained as their actual fields. Empty content wrappers add neither a primitive actor nor a material response. In particular no conclusion about actual receiver response availability follows solely from equal leaf lists.

\paragraph{No controller quotient.}
None of these assertions identifies complete controllers across multiplication. Two histories whose latent grouping is erased can carry different later commands or whole-law choices. The geometry and inverse-image response construction do not supply descent for those tables. Initial-law labels, selected full command laws, zero histories and unselected histories remain in the separate complete-control interface.

\section{Global incidence is additional data, not a fibrewise consequence}
The reserve construction at one event/boundary does not by itself supply a global protocol. Here is an explicit sufficient assembly. Fix a small law-free allocation-boundary graph $H$: its vertices $h$ retain world and allocation-history index, and an atomic edge $e:h\to h'$ retains the run-qualified event, physical boundary and paid/pure-input kind. The actual predecessor relation on allocation indices makes names born at $e$ fresh there and old at subsequent events. The index counts declared atoms, not additional paid events. Free paths in $H$ retain every such atom. Taking $H$ to be the graph of finite declared boundary histories is one concrete choice; no probabilities are assigned to histories.

A global intrinsic protocol supplies COMMON vertex sorts $V_h$ and common entry-symbol sorts $E_h$ with $\iota:E_h\to V_h$, their operational fields and closed marks. For each $e:h\to h'$ it supplies intrinsic atomic records $A_e$, positive ticket requirements and actual incidence
\[
                       V_h\ \longleftarrow A_e\ \longrightarrow V_{h'}.
\]
Entries referenced by a record are actual elements of the SAME $E_{h'}$, not independently selected copies in each event fibre. The source/target, entry, support and observation fields used in the local reserve construction are restrictions of these common fields. Requirements mention only names old at $e$; new names are separate fields stamped at $e$. A ticketed edge is $(t,a)$ with $t\models r_e(a)$. Its endpoints are the common vertices just specified. Local tagged endpoint representatives forget their tag to these same vertices. Any permitted vertex/entry eligibility requirements at $e$ are included in $r_e(a)$ and the associated local incidence requirements, not used to assign contradictory global fields.

A global map consists of ONE map on every $V_h,E_h$, and maps on every $A_e$, commuting with these incidences and with all local weakening maps. Require the previous exact operational and closed-mark fields and ticket-eligibility implications at every $e$. This is an independently defined category of many-sorted diagrams and maps, not a collection of unrelated eventwise isomorphisms.

\begin{proposition}[Global graph assembly]
The atomic pending polynomial, unit and multiplication extend to this global category. They preserve its strict pullbacks and cartesian naturality squares. They induce coherent maps on its free paths; no identification of a path with one atomic edge is made.
\end{proposition}
\begin{proof}
Define vertices and entries at each $h$ by the same support-separated tuples, retaining $h$ even when empty. Define atomic edges at $e$ by the actual common-ticket component edges and designated target entries. Its source and target are exactly the tuples of their common source and target vertices. Thus the local constructions commute with global incidence, including entries already exposing ports. Uniform maps act through the common component maps. The previously proved unit and flattening formulas on vertices and entries are independent of $e$; each local atomic formula has precisely those sources and targets. Consequently they are global graph maps, not just fibrewise maps with potentially mismatched endpoints. Their equations follow on every sort. Pullbacks are pairs with equal images on these same sorts; the local unique reconstructions have the supplied common endpoints, so the cartesian proofs extend. A graph map commuting with endpoints extends uniquely to the free path map, preserving the entire event word.\end{proof}

Independent eventwise data without this common incidence do not imply the proposition. For example the terminal endpoint of one fibre and source endpoint of another can have the same printed value but different primitive supports or different closed marks. Identifying those endpoints would violate strict typing; retaining them separately permits no composable path. This is a genuine gluing obligation. Conversely a receiver used with the global construction must supply its initial and all-raw post-observations into these common sorts and must respect the allocation/history index. A receiver presented only by fixed-boundary observations has not thereby supplied them. Adding an allocation-history presentation is permitted as declared interface enrichment, not claimed to be an isomorphism of opaque receiver states. No law, controller descent or response along a missing path follows from global graph closure.

\end{document}
